\documentclass[10pt,twocolumn,a4paper]{article}

\usepackage[utf8]{inputenc}
\usepackage[margin=0.75in]{geometry}
\usepackage{amsmath,amssymb,amsfonts}
\usepackage{booktabs}
\usepackage{cite}
\usepackage{hyperref}
\usepackage{xcolor}
\usepackage{enumitem}
\usepackage{titlesec}
\usepackage{microtype}

\hypersetup{
    colorlinks=true,
    linkcolor=blue!70!black,
    citecolor=blue!70!black,
    urlcolor=blue!70!black
}

\titleformat{\section}{\large\bfseries}{\thesection}{1em}{}
\titleformat{\subsection}{\normalsize\bfseries}{\thesubsection}{1em}{}

\setlist{noitemsep,topsep=2pt,parsep=2pt,partopsep=2pt}

\title{\textbf{Project BankShield: Stress-Testing Banking AML and Fraud Rules Through Multi-Agent Simulation}\\[0.3em]
\large An Architecture for Finding Policy Blind Spots in Transaction Monitoring}
\author{\textbf{Candidate Brief} \\ Technical Assessment Submission \\ Code repository: \url{https://github.com/Charukhesh/BankShield.git}}
\date{\today}

\begin{document}

\maketitle

\begin{abstract}
Financial criminals are starting to use automated agents and synthetic profiles to probe transaction monitoring systems, finding routes that slip past fixed thresholds. BankShield is an adversarial testing environment designed to discover these gaps proactively. It pairs an exploratory language model agent that drafts multi-hop transaction graphs with a deterministic ledger and rule evaluator that checks account balances and statutory thresholds. This separation keeps the evaluation reproducible and auditable for bank compliance teams, avoiding the arithmetic and logic errors common in pure model pipelines.
\end{abstract}

\section{Problem statement and domain context}

\subsection{The area explored and initial findings}
Retail and commercial banks must monitor payments under the Bank Secrecy Act (31 U.S.C. \S 5318) and Financial Action Task Force standards to detect money laundering, fraud, and structured deposits. 

In reviewing industry data, the starkest finding is the operational imbalance of current systems. The Financial Action Task Force reported in their review of supervisory technology that over 95\% of alerts generated by traditional rules-based and supervised machine learning monitoring systems turn out to be false positives \cite{fatf2021tech}. 

Current rule maintenance is almost entirely reactive. Compliance teams update detection thresholds after law enforcement publishes an advisory or following a regulatory audit penalty. Meanwhile, bad actors can run automated scripts to test transaction limits, splitting funds into sub-\$10,000 transfers across newly created accounts to avoid triggering a Currency Transaction Report.

\subsection{Users, operational tasks, and current costs}
The primary users are Financial Crime Compliance officers, transaction monitoring analysts, and model risk validation teams. 

Their recurring task is calibrating and validating detection rules against new laundering typologies. In most banks, this testing occurs quarterly or annually through manual retrospective sample reviews.

The financial cost of this workflow is substantial. LexisNexis Risk Solutions estimated that financial crime compliance cost global institutions over \$206 billion in 2023 \cite{lexisnexis2023}. Most of that spending goes toward large operational teams manually clearing false alerts. Meanwhile, penalties for undetected laundering failures can reach hundreds of millions of dollars.

\subsection{Scope boundaries}
\begin{itemize}
    \item \textbf{In scope:} Payment business logic, structuring heuristics, transfer velocity limits, synthetic account graphs, and automated gap reporting.
    \item \textbf{Out of scope:} Network-level security, operating system vulnerabilities, cryptographic protocol attacks, and physical branch or ATM controls.
\end{itemize}

\section{Solution design and systems architecture}

\subsection{System overview}
BankShield splits the testing pipeline into three parts:
\begin{enumerate}
    \item An adversarial scenario generator using language models to construct structured transaction graphs (such as smurfing rings and multi-hop layering routes).
    \item A deterministic state engine that runs these transactions against real bank rules and double-entry accounting balances.
    \item An audit reporting module that maps confirmed gaps to specific internal rule IDs and regulator guidelines.
\end{enumerate}

\subsection{Separation of probabilistic and deterministic tasks}
A primary design choice in BankShield is that language models do not decide whether a rule was broken. Models make arithmetic mistakes and drift across runs, which fails bank auditability standards.

The language model handles the probabilistic task: inventing plausible counterparty networks, timing variations, transaction descriptions, and amounts designed to avoid simple flags. 

Deterministic Python code handles the mathematical verification. It checks explicit equations, such as whether cumulative deposits to one beneficiary exceed \$10,000 within 48 hours, or whether a routing path touches high-risk foreign institutions. Every test run uses fixed random seeds so auditors can reproduce the sequence step by step.

\subsection{Knowledge storage and regulatory updates}
Statutory rules and compliance thresholds live directly in version-controlled configuration files. Public advisories from FinCEN and the FATF are indexed as reference text for the reporting module.

When regulators change a requirement, engineers update the corresponding schema in the rule registry. The scenario generator then uses the updated constraints on its next run.

Outputs are recorded as structured JSON graphs detailing account identifiers, timestamps, transfer amounts, and the rule conditions that failed to catch them.

\subsection{System architecture}
\begin{center}
\fbox{\parbox{0.92\linewidth}{
\centering \small
\textbf{Adversarial Scenario Generator (LLM)} \\
Generates synthetic accounts, amounts, and transfer routes \\
$\downarrow$ \\
\textbf{Deterministic Bank Rule Engine} \\
Executes math, ledger balances, and threshold checks \\
$\downarrow$ \\
\textit{Was an evasion path found?} \\
$\rightarrow$ \textbf{No:} Rules caught the attack. \\
$\rightarrow$ \textbf{Yes:} Unflagged transfer sequence recorded. \\
$\downarrow$ \\
\textbf{Audit Reporter Module} \\
Links the gap to specific rule IDs and FinCEN advisories \\
$\downarrow$ \\
\textbf{Compliance Officer Review} \\
Reviews diff and updates production thresholds
}}
\end{center}

\subsection{Alternatives considered and rejected}
We evaluated three other approaches before settling on this architecture:
\begin{itemize}
    \item \textbf{End-to-end evaluation inside a language model:} We rejected this because models miscalculate rolling sums across dozens of synthetic accounts and produce inconsistent decisions on the same transaction log.
    \item \textbf{Random parameter fuzzing:} Traditional fuzzers mutate inputs randomly, which mostly creates invalid data (like negative amounts or malformed account IDs) that standard input validation catches immediately. They rarely stumble into coordinated, multi-step laundering structures.
    \item \textbf{Supervised machine learning models:} Supervised models rely on past labeled data. They flag known patterns well, but they offer little insight into novel evasion routes.
\end{itemize}

\section{Risk analysis and mitigation}

\begin{table*}[t]
\centering
\small
\begin{tabular}{@{}p{3.0cm}p{4.4cm}p{9.4cm}@{}}
\toprule
\textbf{Risk} & \textbf{Failure Mode} & \textbf{Engineering Mitigation} \\ \midrule
Erroneous findings & The model suggests a scenario that violates accounting laws. & The system routes every generated transfer through a double-entry ledger. If a sequence fails balance conservation ($Assets = Liabilities + Equity$), it is dropped immediately. \\ \midrule
Outdated rules & Testing against superseded regulatory guidelines. & Rule definitions are stored in Git with release tags tied to legal enforcement dates. \\ \midrule
Confidential data & Exposing private customer payment histories. & BankShield uses procedural synthetic data generated from statistical models. It does not ingest production customer records or personal identifiers. \\ \midrule
Dual-use risk & Attack vectors could be used by actual criminals. & The platform runs in isolated test environments. Reports output rule patches and coverage analyses rather than operational exploit playbooks. \\ \midrule
Regulatory liability & Erroneous automated changes to production compliance rules. & BankShield operates in an advisory mode. It generates recommended rule patches that require review and approval from a designated compliance officer. \\ \bottomrule
\end{tabular}
\caption{System risks and technical controls.}
\label{tab:risks}
\end{table*}

\subsection{Questions to resolve before production use}
Two practical questions need verification before expanding this system:
First, how well does this testing fit existing model validation guidelines? In the United States, Federal Reserve Supervisory Letter SR 11-7 sets strict requirements for conceptual soundness and ongoing monitoring \cite{fedsr117}. Bank compliance teams must determine whether synthetic adversarial testing satisfies their internal model validation policies.

Second, does the system generate economically realistic attacks? An agent might discover that splitting \$10,000 into 10,000 one-dollar transactions bypasses a velocity rule. In practice, bank transfer fees make that attack impractical for launderers. The test generator must factor in real-world friction so compliance teams do not waste time on implausible edge cases.

\newpage

\section*{Appendix: What we added beyond the AI}
\addcontentsline{toc}{section}{Appendix: What we added beyond the AI}

\subsection*{Where AI helped}
We used language models to brainstorm adversarial laundering setups and write synthetic merchant names and payment memos. Two prompts that shaped this were:
\begin{itemize}
    \item ``List the common structural patterns money laundering networks use to move funds across multiple personal accounts while staying below the \$10,000 Currency Transaction Report threshold.''
    \item ``Draft a realistic sequence of transactions between eight synthetic consumer accounts, including normal household expenses alongside structured transfers.''
\end{itemize}

\subsection*{What we rejected or corrected}
The language model originally suggested letting an LLM review the final transaction logs and decide which rules were violated. We rejected that recommendation because language models cannot reliably calculate running sums across multiple accounts and currencies. We replaced that step with a deterministic Python rules engine.

The model also suggested writing an automated patch system to alter bank rule sets in real time. We discarded that proposal. Banking regulations require human approval and documented model risk reviews before any production rule change.

Finally, early prompt responses generated detailed instructions on how to evade specific bank detection filters. We added system prompts and input constraints to ensure the generator only produces synthetic test fixtures for compliance verification.

\subsection*{Decisions from domain research}
\begin{itemize}
    \item Enforcing double-entry ledger bookkeeping on all generated transactions to guarantee that synthetic tests obey basic mathematical and accounting laws.
    \item Focusing on small and mid-sized community banks, which often bear the highest relative compliance costs because they lack dedicated internal red teams.
    \item Building explicit JSON DAG schemas for all test outputs so findings can be handed directly to internal auditors without reformatting.
\end{itemize}

\subsection*{One element defensible without AI}
The core architecture: separating scenario generation from a deterministic, double-entry state engine that checks statutory BSA and AML rules. Even if the scenarios were generated by rule permutations or manual test scripts rather than an LLM, this decoupled structure gives compliance teams an auditable, reproducible test harness.

\subsection*{Code repository}
The source code, rule engine implementation, and synthetic test suites are publicly accessible on GitHub at \url{https://github.com/Charukhesh/BankShield.git}.

\begin{thebibliography}{99}
\bibitem{ustreasury2024}
U.S. Department of the Treasury, ``Managing Artificial Intelligence-Specific Cybersecurity Risks in the Financial Services Sector,'' Office of Cybersecurity and Critical Infrastructure Protection, Mar. 2024.

\bibitem{fatf2021tech}
Financial Action Task Force, ``Opportunities and Challenges of New Technologies for AML/CFT,'' Paris, France, 2021.

\bibitem{lexisnexis2023}
LexisNexis Risk Solutions, ``True Cost of Financial Crime Compliance Global Report,'' 2023.

\bibitem{fincen2021}
Financial Crimes Enforcement Network, ``Advisory on Environmental Crimes and Associated Illicit Financial Activity,'' FIN-2021-A003, Nov. 2021.

\bibitem{weber2019graph}
M. Weber et al., ``Anti-Money Laundering in Bitcoin: Experimenting with Graph Convolutional Networks for Financial Forensics,'' \textit{arXiv preprint arXiv:1908.02591}, 2019.

\bibitem{fedsr117}
Board of Governors of the Federal Reserve System, ``Supervisory Guidance on Model Risk Management,'' SR Letter 11-7, 2011.
\end{thebibliography}

\end{document}
